\section{Integral orders and the imported block constructions}\label{sec:inputs}

The proof requires numerical bounds for actual blocks and a precise
description of the crossed orders produced by reduction.  We state
these inputs separately.  The distinction will matter when the
numerical criterion is applied: first we identify a restriction as a
group block, and then we use its Cartan and Frobenius bounds.

\subsection{Orders, Frobenius and the numerical criterion}

An $\cO$-order means a unital $\cO$-algebra free of finite rank as an
$\cO$-module.  For a general order, a block summand means its direct
factor at a primitive central idempotent.  Put $K_0=\Frac(\cO)$ and
let $\sigma$ be Witt
Frobenius.  For an order $A$, the coefficient twist $\sigma^m A$ is
the order obtained by applying $\sigma^m$ to its structure constants;
equivalently it is scalar transport along $\sigma^m$.  For a group
block we identify it with $\cO G\sigma^m(b)$.  We write
$\mf_{\cO}(A)$ for the least positive $m$ for which $A$ and
$\sigma^m A$ are $\cO$-linearly Morita equivalent, whenever such an
$m$ exists.  Only upper bounds for this number will be used.

Every finite $\cO$-algebra is semiperfect.  A basic idempotent of a
block order is a sum of one primitive idempotent for every isomorphism
type of indecomposable projective.  It is full, and its corner is the
basic order.  Basic orders are unique up to isomorphism within a
Morita class.  If the Cartan matrix of a block is $(c_{ij})$, its
basic order has $\cO$-rank $\sum_{i,j}c_{ij}$: reduction is a split
basic $k$-algebra, whose dimension is that sum.

\begin{proposition}[Cartan and integral finiteness inputs]\label{prop:criteria}
The following assertions will be used.
\begin{enumerate}
\item For every $p$ and $M$, the Cartan sums of all finite-group
blocks over $\cO$ with defect order at most $M$ are bounded by a
constant $c(p,M)$.
\item For fixed positive bounds on the defect order, Cartan sum,
and integral Morita--Frobenius number, finite-group blocks over
$\cO$ have only finitely many $\cO$-linear Morita equivalence
classes.
\end{enumerate}
\end{proposition}

\begin{proof}
By \cite[Theorem~1.1]{OpenAIFieldDonovan2026}, the reductions of the
blocks in (i) have finitely many $k$-Morita classes.  Their Cartan
matrices, up to simultaneous permutation, consequently form a finite
list.  Integral blocks and their reductions have the same Cartan
matrix, which proves (i).

Assertion (ii) is \cite[Theorem~3.10]{EEL2020}, applied to the
finitely many possible defect exponents.  The coefficient conventions
of that theorem include the absolutely unramified complete discrete
valuation ring $W(k)$ and its coefficient Frobenius, which fixes the
uniformizer $p$.
Its defect-zero case can also be separated: a defect-zero block over
$W(k)$ is a matrix algebra over $W(k)$ and has basic order $\cO$.
Indeed, its reduction is a full matrix algebra over $k$; lift its
matrix units over the complete ring $\cO$, whose resulting primitive
corner has rank one.
\end{proof}

\subsection{Multiplication factors and identity markings}

\begin{definition}[Crossed orders and based isomorphisms]
Let $R$ be an $\cO$-order and $Q$ a finite group.  A normalized
crossed system consists of $\alpha_q\in\Aut_{\cO}(R)$ and
$a_{q,t}\in R^\times$ such that
\begin{align}
 \alpha_q\alpha_t&=\ad(a_{q,t})\alpha_{qt},
       \label{eq:action-law}\\
 a_{q,t}a_{qt,v}&=\alpha_q(a_{t,v})a_{q,tv},
       \label{eq:factor-law}
\end{align}
with $\alpha_1=\id$ and $a_{1,q}=a_{q,1}=1$.
Its crossed order is $A=\bigoplus_{q\in Q}Ru_q$, with
\[
 u_qr=\alpha_q(r)u_q,\qquad
 u_qu_t=a_{q,t}u_{qt},\qquad u_1=1.
\]
A $Q$-labelled graded isomorphism preserves each homogeneous
component with its label.  When the identity component is identified
with a fixed $R$, such an isomorphism is \emph{based} if it is the
identity on $R$.
\end{definition}

Changing each $u_q$ to $t_qu_q$, with $t_q\in R^\times$ and
$t_1=1$, gives exactly the based changes of crossed systems.  The
induced homomorphism $Q\to\Out_{\cO}(R)$ is the outer action.
It remembers the actions only modulo inner automorphisms.  Both the
units implementing their products in Equation~\eqref{eq:action-law}
and the compatibility in Equation~\eqref{eq:factor-law} are needed to
recover the order.  For $S\leq Q$ the \emph{restriction to $S$} is
the crossed suborder $A_S=\bigoplus_{s\in S}Ru_s$, with the same
identity marking and the restricted factors.

\subsection{Reduction and dual recovery}

We call a pair $(H,B)$ reduced if $B$ is a block of $\cO H$ such
that every block of a normal subgroup covered by $B$ is $H$-stable,
and
\begin{equation}\label{eq:reduced}
 B\text{ covers a nilpotent block of }H_0\trianglelefteq H
 \quad\Longrightarrow\quad H_0\leq Z(H)O_p(H).
\end{equation}
An--Eaton's preliminary reduction
\cite[Proposition~6.1]{AnEaton2025} replaces an arbitrary block by a
reduced pair through an $\cO$-linear basic Morita equivalence,
preserving the defect group.  Their coefficient conventions include
$\cO=W(k)$.  We apply this integral reduction directly.

The next proposition collects the group structure, actual factor
elements and integral full corners supplied by the field companion.
Constants in it depend only on $p,M$.

\begin{proposition}[Controlled integral extensions]\label{prop:structure}
Let $(H,B)$ be a reduced pair with defect order at most $M$.  Put
$N=F^*(H)$, $X=H/N$, and let $b$ be the block of $\cO N$ covered by
$B$.  Then the following constructions and bounds hold.
\begin{enumerate}
\item The group $N$ is a central product
\[
 N=PZK_1\cdots K_j,\qquad P=O_p(H),\quad Z=O_{p'}(H)\leq Z(H),
\]
where the $K_i$ are quasisimple components, $|P|$ and $j$ are
bounded, and $[X:O_{p'}(X)]$ is bounded.  The latter index is
bounded also for every subgroup of every extension of $X$ by a
$p'$-group.
\item There is an extension $N\trianglelefteq\bar N$ with abelian
$p'$-quotient $A_0=\bar N/N$.  For representatives $h_x$ of $x\in X$,
normalized by $h_1=1$, write $h_xh_y=n_{x,y}h_{xy}$ with
$n_{x,y}\in N$.  Conjugation extends to automorphisms $\alpha_x$ of
$\bar N$ satisfying the actual identities
\begin{align}
 \alpha_x\alpha_y&=\ad(n_{x,y})\alpha_{xy},
       \label{eq:actual-actions}\\
 n_{x,y}n_{xy,z}&=\alpha_x(n_{y,z})n_{x,yz}.
       \label{eq:actual-factors}
\end{align}
Every block $\bar b$ of $\cO\bar N$ covering $b$ has bounded defect.
\item Put $\Xi=\Hom(A_0,k^\times)$, using Teichm\"uller lifts of
the characters over $\cO$, and $Y=\Xi\rtimes X$.  There is a
$Y$-crossed order with identity component $\cO\bar N$ in which
$\cO H$ is a full idempotent corner.  After taking an appropriate
central summand and a further full corner, this is a crossed order
with identity component
\[
 B_0=\cO\bar N\bar b
\]
and grading group $Q=\Stab_Y(\bar b)$.  In the pure $X$-degrees the
actions and factors are those in (ii).  The index
$[Q:O_{p'}(Q)]$ is bounded.
\item The basic orders of all these identity blocks $B_0$ belong
to a finite list of integral orders $R_1,\ldots,R_t$.  Their integral
Picard groups, and in particular their outer automorphism groups,
are finite.
\end{enumerate}
\end{proposition}

\begin{proof}[Source of the assertions]
For a reduced pair, the structural conclusions follow from
\cite[Lemma~8.1]{OpenAIFieldDonovan2026}.  Its opening Morita
reduction is stated over $k$; we use only its conclusions about the
already reduced pair, obtained here by the integral An--Eaton
reduction.  The compatible partial extension and its actual factors
are \cite[Lemma~8.2]{OpenAIFieldDonovan2026}.  The integral dual
recovery construction and full corners are
\cite[Lemma~8.3]{OpenAIFieldDonovan2026}.

For clarity, the dual-recovery idempotent is
$|\Xi|^{-1}\sum_{\chi\in\Xi}u_\chi$.  The characters are trivial
on every $n_{x,y}\in N$, so the pure $X$-units preserve this
idempotent and retain exactly these factors.  This explains why
the construction is integral and why it retains the multiplication
data needed later.  Finally (iv) is the integral assertion of
\cite[Lemma~8.5]{OpenAIFieldDonovan2026}, using the trivial-subgroup
case of its Proposition~9.1.  The finite Picard assertion is also
\cite[Theorem~B and Corollary~1.2]{EiseleGeometry2022}.
\end{proof}

Every $p$-subgroup $S$ of $Y$ can be conjugated by an element of
$\Xi$ into the pure subgroup $X$.  Indeed, its image in $X$ is a
$p$-group, and Schur--Zassenhaus conjugates the two complements to
$\Xi$ in the inverse image of that group.  In the crossed order this
conjugation is implemented by an integral homogeneous unit; it also
transports $\bar b$.  We may therefore work with pure $p$-subgroups,
provided we keep the transported identity block in the same family.
In particular, the orders of the relevant $p$-subgroups of $Q$ are
bounded: their intersection with $O_{p'}(Q)$ is trivial, so they
inject into the quotient whose order was bounded in (iii).

There is no bound here on the orders of the central kernels in the
universal-cover presentations, or on the extra central $p$-factors
used in the partial regular extensions.  The comparison in
Section~\ref{sec:comparison} descends through those kernels before
any bounded-defect integral criterion is applied.

\subsection{The companion's comparison data}

The comparison convention in
\cite[Definition~8.4 and Proposition~9.1]{OpenAIFieldDonovan2026}
uses an integral Morita bimodule but asserts exact covariance and
factor identities on its reduction.  We use the particular integral
operators constructed there.  Their ingredients have four distinct
roles.  Lemmas~9.2--9.4 realize the genuine field, graph and inner
relations, construct an invariant torus and character, and choose a
common parabolic--Levi pair with the specified inner overlap.
Lemma~9.5 gives an integral Levi Morita bimodule with strict geometric
actions and matching central actions.  Lemma~9.6 gives a uniformly
bounded coefficient-Frobenius return and its prime-to-$p$ character
twist.  Finally, Lemma~9.7 computes the integral scalar overlap of
the chosen operators.

Section~\ref{sec:comparison} recalls these constructions with their
exact hypotheses and shows that their scalar errors stay in $\T$.
The subsequent normalization over $\cO$ is proved here.  Thus the
imported comparison supplies operators and overlap identities, while
Proposition~\ref{prop:comparison} supplies the exact integral factor
law required by the later extension argument.
